\chapter{Fungsi Peubah Real}\label{ch:b2:realfun}

Sebelum analisis beralih ke fungsi atas fungsi (\cref{ch:b2:funcseq}),
ada gunanya mengenal lanskap satu peubah lebih rinci daripada yang
dituntut Tahun ke-1: seberapa tidak \omterm{def:b2:metric:continuity}{kontinu} sebuah fungsi monoton bisa
jadi, seberapa teratur sebuah fungsi cembung harus jadi, dan sifat
istimewa apa yang dinikmati turunan (Darboux). Hasil struktural ini
singkat, tajam, dan disukai para penguji.

\section{Fungsi monoton}

\begin{theorem}[Keteraturan fungsi monoton]\label{thm:b2:realfun:monotone}
Misalkan $f \colon I \to \R$ naik pada sebuah interval.
\begin{enumerate}
  \item Di setiap titik dalam $a$, limit sepihaknya ada:
        \[
        f(a^-) = \sup_{x < a} f(x) \;\leq\; f(a) \;\leq\;
        f(a^+) = \inf_{x > a} f(x) ;
        \]
        dan setiap ketakkontinuannya berupa \emph{lompatan}.
  \item Himpunan ketakkontinuan $f$ paling banyak
        \omterm{def:b2:structures:countable}{terbilang}.\index{fungsi monoton!ketakkontinuan}
\end{enumerate}
\end{theorem}

\begin{proof}
(1) Himpunan $\{f(x) : x < a\}$ tak kosong dan terbatas di atas oleh
$f(a)$: supremumnya $s$ memenuhi $f(x) \to s$ ketika $x \to a^-$
(diberikan $\varepsilon$, ada $f(x_0) > s - \varepsilon$, lalu
kemonotonannya menjebak $f(x) \in \intoc{s - \varepsilon}{s}$ untuk $x
\in \intoo{x_0}{a}$). Setangkup dengan itu di sebelah kanan.

(2) Kepada tiap ketakkontinuan $a$ lekatkan interval \omterm{def:b2:metric:topology}{terbuka} tak kosong
$J_a = \intoo{f(a^-)}{f(a^+)}$ (yakni lompatan yang sejati). Untuk dua
ketakkontinuan $a < b$, himpunan $J_a$ dan $J_b$ saling lepas: sebab
$f(a^+) \leq f(c) \leq f(b^-)$ untuk sembarang $c$ di antaranya. Tiap
$J_a$ memuat sebuah bilangan rasional, dan ketakkontinuan yang berbeda
mendapat rasional yang berbeda: jadi ada injeksi himpunan
ketakkontinuannya ke $\Q$, yang \omterm{def:b2:structures:countable}{terbilang}
(\cref{prop:b2:structures:countablestable}).
\end{proof}

\begin{example}
Batasnya tajam: tetapkan sebuah pencacahan $(r_n)$ bagi $\Q \cap
\intoo{0}{1}$ lalu tetapkan $f(x) = \sum_{n : r_n \leq x} 2^{-n}$
(yakni definisi lewat \omterm{def:b2:series:summable}{keluarga terjumlahkan},
\cref{def:b2:series:summable}). Maka $f$ naik pada
$\intcc{0}{1}$ dan tidak \omterm{def:b2:metric:continuity}{kontinu} persis di setiap bilangan rasional
$\intoo{0}{1}$ (dengan lompatan $2^{-n}$ di $r_n$): jadi fungsi monoton
\emph{dapat} tidak \omterm{def:b2:metric:continuity}{kontinu} pada \omterm{def:b2:structures:countable}{himpunan terbilang} yang padat.
\end{example}

\begin{example}[Lompatannya tak dapat melebihi kenaikannya]\label{ex:b2:realfun:jumpbudget}
Untuk $f$ yang naik pada $\intcc{a}{b}$, lompatannya punya anggaran:
jika $a < c_1 < \dots < c_m < b$ merupakan ketakkontinuan dengan lompatan
$s_i = f(c_i^+) - f(c_i^-) > 0$, maka dengan memilih titik
selang-seling $a < c_1 < t_1 < c_2 < \dots$ lalu memakai kemonotonan
pada tiap potongannya,
\[
\sum_{i=1}^{m} s_i \;\leq\; f(b) - f(a) :
\]
jadi kenaikan totalnya membatasi lompatan totalnya. Akibatnya: untuk
tiap $k$, paling banyak $k\,\bigl(f(b) - f(a)\bigr)$ ketakkontinuan
berlompatan $\geq \frac1k$ --- yakni perhalusan kuantitatif bagi
\cref{thm:b2:realfun:monotone}~(2), sebab himpunan ketakkontinuannya
adalah gabungan \omterm{def:b2:structures:countable}{terbilang} atas $k$ dari himpunan hingga itu. Pada
fungsi berlompatan rasional di atas, anggarannya terpakai habis:
lompatan $2^{-n}$ berjumlah $1 = f(1^+) - f(0^-)$ dalam arti perluasan
yang jelas. Jadi fungsi monoton boleh melompat secara padat, tetapi hanya
dengan tunjangan yang ketat.
\end{example}

\section{Fungsi cembung}

\begin{lemma}[Ketaksamaan kemiringan]\label{lem:b2:realfun:slopes}
Misalkan $f$ cembung pada $I$ dan $x < y < z$ di $I$. Maka
\[
\frac{f(y) - f(x)}{y - x}
\;\leq\; \frac{f(z) - f(x)}{z - x}
\;\leq\; \frac{f(z) - f(y)}{z - y} :
\]
jadi kemiringan tali busurnya naik pada kedua ujungnya.\index{fungsi
cembung!ketaksamaan kemiringan}
\end{lemma}

\begin{proof}
Tulis $y = \frac{z - y}{z - x}\,x + \frac{y - x}{z - x}\,z$: yakni
kombinasi cembung, sebab kedua koefisiennya positif dan
berjumlah $1$. Kecembungannya memberi
\[
f(y) \;\leq\; \frac{z-y}{z-x}\,f(x) + \frac{y-x}{z-x}\,f(z).
\]
Untuk ketaksamaan kirinya, kurangkan $f(x)$ dari kedua ruasnya, dengan
memakai $\frac{z-y}{z-x} - 1 = -\frac{y-x}{z-x}$:
\[
f(y) - f(x) \leq \frac{y - x}{z - x}\bigl(f(z) - f(x)\bigr),
\]
lalu bagilah dengan $y - x > 0$. Untuk ketaksamaan kanannya, kurangkan
justru dari $f(z)$:
\[
f(z) - f(y) \geq f(z) - \frac{z-y}{z-x}f(x) -
\frac{y-x}{z-x}f(z)
= \frac{z - y}{z - x}\bigl(f(z) - f(x)\bigr),
\]
lalu bagilah dengan $z - y > 0$. Kedua langkah yang tertulis itu adalah
kesamaan barisentrik yang sama, dibaca terhadap ujung yang berbeda.
\end{proof}

\begin{theorem}[Keteraturan fungsi cembung]\label{thm:b2:realfun:convexreg}
Misalkan $f$ cembung pada sebuah interval $I$.
\begin{enumerate}
  \item Di setiap titik dalam, $f$ punya turunan sepihak yang
        hingga $f'_g \leq f'_d$; keduanya merupakan fungsi naik
        terhadap titiknya; khususnya $f$ \omterm{def:b2:metric:continuity}{kontinu} pada bagian
        dalam $I$ (tetapi boleh jadi tidak di ujungnya).
  \item Fungsi $f$ terletak di atas tiap \emph{garis penyangga}nya:
        untuk $a$ yang di dalam dan sembarang $m \in
        \intcc{f'_g(a)}{f'_d(a)}$,
        \[
        f(x) \geq f(a) + m(x - a) \qquad (x \in I).
        \]
  \item (Jensen, berbobot) Untuk $x_i \in I$ dan bobot $\lambda_i
        \geq 0$ dengan $\sum\lambda_i = 1$:
        \[
        f\Bigl(\sum_i \lambda_i x_i\Bigr) \leq \sum_i \lambda_i
        f(x_i) .\index{ketaksamaan Jensen}
        \]
\end{enumerate}
\end{theorem}

\begin{proof}
(1) Tetapkan $a$ di bagian dalam. Menurut \cref{lem:b2:realfun:slopes},
kemiringan $\tau(h) = \frac{f(a + h) - f(a)}{h}$ merupakan fungsi naik
terhadap $h$ (pada kedua sisinya, dan $\tau(h_-) \leq \tau(h_+)$ untuk
$h_- < 0 < h_+$). Karenanya $\tau$ punya limit hingga ketika $h \to 0^-$
(sebab naik dan terbatas di atas oleh kemiringan kanan mana pun) ---
itulah $f'_g(a)$ --- dan juga ketika $h \to 0^+$ (itulah $f'_d(a)$),
dengan $f'_g(a) \leq f'_d(a)$. Turunan sepihak yang hingga memaksa
kekontinuan di $a$. Kemonotonan terhadap titiknya: untuk $a < b$ yang di
dalam, $f'_d(a) \leq \frac{f(b) - f(a)}{b - a} \leq f'_g(b)$, lagi-lagi
lewat ketaksamaan kemiringannya.

(2) Untuk $x > a$: $\frac{f(x) - f(a)}{x - a} \geq f'_d(a) \geq m$;
sedangkan untuk $x < a$: $\frac{f(a) - f(x)}{a - x} \leq f'_g(a) \leq
m$. Keduanya tertata ulang menjadi klaimnya.

(3) Induksi pada banyaknya titik, persis seperti pada jilid Tahun ke-1
(sebab kasus dua titiknya adalah definisinya) --- atau dalam satu
tarikan: terapkan (2) di $a = \sum\lambda_i x_i$ lalu rata-ratakan
ketaksamaan garis penyangganya di titik $x_i$ dengan bobot $\lambda_i$:
$\sum_i \lambda_i f(x_i) \geq f(a) + m\sum_i\lambda_i(x_i - a) = f(a)$.
\end{proof}

\begin{omfigure}
\begin{tikzpicture}
\begin{axis}[omaxis, width=8.2cm, height=5.4cm,
    xmin=-2.4, xmax=2.5, ymin=-0.9, ymax=4.5,
    xtick={-1.5, 0.5, 2}, ytick=\empty]
  \addplot[omDef, very thick, domain=-2.1:2.15, samples=80]
    {x^2};
  \addplot[omThm, thick, domain=-1.5:2, samples=2]
    {0.5*x + 3};
  \addplot[omProp, thick, domain=-1.1:2.2, samples=2]
    {x - 0.25};
  \node[font=\small, omThm] at (axis cs:-1.05,3.15)
    {tali busur};
  \node[font=\small, omProp] at (axis cs:1.75,0.95)
    {garis penyangga};
  \node[font=\small, omDef] at (axis cs:-1.95,2.6) {$f$};
\end{axis}
\end{tikzpicture}

{\small Kecembungan dalam satu gambar: di antara $-1.5$ dan $2$,
grafik $f(x) = x^2$ tetap di bawah tali busurnya (yakni definisinya) dan
di atas garis penyangganya di $x = 0.5$
(\cref{thm:b2:realfun:convexreg}~(2)) --- setiap ketaksamaan pada soal
akhir pekan bab ini tak lain penataan ulang kedua kedudukan
itu.}
\end{omfigure}

\begin{example}[Ketakkontinuan di ujung]
Pada $\intcc{0}{1}$, fungsi dengan $f(0) = 1$ dan $f(x) = 0$ untuk $x >
0$ bersifat cembung tetapi tidak \omterm{def:b2:metric:continuity}{kontinu} di ujung $0$: jadi pernyataan
(1) memang tajam.
\end{example}

\begin{example}[Sudut dan berkas garis penyangga]\label{ex:b2:realfun:corner}
Untuk $f(x) = \abs x$ di $a = 0$: turunan sepihaknya adalah
$f'_g(0) = -1$ dan $f'_d(0) = +1$, lalu
\cref{thm:b2:realfun:convexreg}~(2) membagikan garis penyangga bagi
\emph{setiap} kemiringan $m \in \intcc{-1}{1}$:
\[
\abs x \geq m\,x \qquad (x \in \R,\ -1 \leq m \leq 1),
\]
yang masing-masing menjadi kesamaan persis pada sebuah setengah-garis
atau di $0$. Fungsi cembung dapat diturunkan di $a$ persis ketika
berkasnya runtuh menjadi satu garis ($f'_g(a) = f'_d(a)$); sedangkan
sudut mengusung seinterval garis singgung. Berkas ini adalah benih
berdimensi hingga bagi \emph{subdiferensial} pada optimasi cembung ---
sekaligus alasan mengapa fungsi cembung begitu tangguh: sebab walaupun
turunannya gagal, geometri penyangganya bertahan, dan hanya itulah yang
dipakai bukti Jensen.
\end{example}

\begin{example}[Ketaksamaan rata-rata pangkat]\label{ex:b2:realfun:powermeans}
Untuk $0 < p < q$ dan $x_i$ yang positif dengan bobot $\lambda_i$
berjumlah $1$, menerapkan Jensen pada $t \mapsto t^{q/p}$ yang cembung
di titik $x_i^p$ memberi
\[
\Bigl(\sum \lambda_i x_i^{p}\Bigr)^{1/p}
\leq \Bigl(\sum \lambda_i x_i^{q}\Bigr)^{1/q} :
\]
jadi rata-rata pangkatnya naik seiring eksponennya --- yang memuat
AM--QM, dan pada limit $p \to 0$ (\cref{exo:b2:realfun:6}) memuat
ketaksamaan AM--GM sekali lagi.
\end{example}

\begin{omfigure}
\begin{tikzpicture}
\begin{axis}[omaxis, width=8.6cm, height=5.6cm,
    xmin=-8.5, xmax=8.5, ymin=0.5, ymax=4.5,
    xtick={-4,-1,1,2,4}, ytick={1,2,4},
    xlabel={$p$}, ylabel={$M_p$}]
  \addplot[omDef, very thick, domain=-8:-0.4, samples=90]
    {((1/3)*(1^x + 2^x + 4^x))^(1/x)};
  \addplot[omDef, very thick, domain=0.4:8, samples=90]
    {((1/3)*(1^x + 2^x + 4^x))^(1/x)};
  \addplot[omThm, dashed, domain=-8.5:8.5, samples=2] {4};
  \addplot[omThm, dashed, domain=-8.5:8.5, samples=2] {1};
  \node[font=\small, omDef] at (axis cs:-5.4,1.55)
    {$M_p(1,2,4)$};
  \node[font=\small, omThm] at (axis cs:-6.4,3.75)
    {$\max = 4$};
  \node[font=\small, omThm] at (axis cs:6.3,1.25)
    {$\min = 1$};
  \node[font=\small, omProp] at (axis cs:0.2,1.83) {$2$};
\end{axis}
\end{tikzpicture}

{\small Rata-rata pangkat $M_p$ atas nilai $1, 2, 4$ (dengan bobot
sama), sebagai fungsi eksponen $p$: naik dari
$\min = 1$ (ketika $p \to -\infty$) ke $\max = 4$ (ketika $p \to
+\infty$), melewati rata-rata harmonik ($p = -1$), geometri (yakni
jurang di $p = 0$, bernilai $2$), aritmetika ($p = 1$) dan kuadratik ($p
= 2$). Seluruh rantai ketaksamaan rata-rata klasik menjadi
satu kurva yang naik --- yang dibuktikan pada soal akhir pekan bab ini,
Bagian III.}
\end{omfigure}

\begin{example}[Entropi maksimal]\label{ex:b2:realfun:entropy}
Untuk vektor peluang $(p_1, \dots, p_n)$ (yang positif dan
berjumlah $1$), entropinya $H(p) = -\sum_i p_i\ln p_i$
memenuhi
\[
H(p) \leq \ln n ,
\qquad\text{dengan kesamaan bila dan hanya bila } p_i = \frac1n \text{ untuk setiap }
i .
\]
Buktinya lewat Jensen (\cref{thm:b2:realfun:convexreg}~(3)) yang
diterapkan pada $\ln$ yang \emph{cekung} dengan bobot $p_i$ di titik
$\frac{1}{p_i}$:
\[
H(p) = \sum_i p_i\ln\frac{1}{p_i}
\leq \ln\Bigl(\sum_i p_i\,\frac1{p_i}\Bigr) = \ln n ,
\]
dan kesamaannya memaksa semua titik $\frac1{p_i}$ sama (lewat
kecekungan tegas), yakni $p$ seragam. Setara dengan itu, inilah
\cref{exo:b2:realfun:7} dengan $q$ yang seragam. Ketakpastiannya
dimaksimumkan oleh ketidaktahuan yang tersebar seragam --- itulah asas
variasional di balik penyandian, mekanika statistik, dan kemunculan
entropi pada \cref{ch:b2:randomvar}.
\end{example}

\begin{method}[Menemukan fungsi cembung di balik sebuah ketaksamaan]\label{met:b2:realfun:spotting}
Sebagian besar ketaksamaan klasik adalah Jensen yang berkostum; untuk
menelanjanginya: (1) normalkan sehingga muncul sebuah \emph{rata-rata
berbobot} (dengan bobot positif berjumlah $1$ --- bagilah dengan massa
totalnya bila perlu); (2) lihatlah fungsi apa yang diterapkan di dalam
dan di luar rata-ratanya: klaim ``$f(\text{rata-rata}) \leq$
rata-rata $f$'' menamai $f$ yang cembung itu; (3) sahkan kecembungannya
lewat turunan keduanya, lalu tangani kesamaannya lewat ketegasannya; (4)
bila tak ada rata-rata yang terlihat, ambillah logaritmanya lebih dulu
--- sebab hasil kali dan pangkat berubah menjadi rata-rata, dan
kecekungan $\ln$ mengusung AM--GM, Young beserta kerabatnya (soal akhir
pekan bab ini menjalankan langkah 1--4 pada masing-masingnya). Jika
logaritma pun tak menyingkap sebuah rata-rata, cobalah membaca
ketaksamaannya sebagai kemonotonan kemiringan
(\cref{lem:b2:realfun:slopes}) --- sebab pernyataan kesuperaditifan
seperti \cref{exo:b2:realfun:9} hidup di sana.
\end{method}

\begin{remark}[Jebakan yang sering muncul]\label{rem:b2:realfun:pitfalls}
(i) Kecembungan tidak awet terhadap hasil kali: $x$ dan $(x - 1)^2$
cembung pada $\intcc{0}{2}$, tetapi hasil kalinya $x(x-1)^2$ punya
turunan kedua $6x - 4$, yang negatif pada
$\intco{0}{\frac23}$ --- jadi tidak cembung; kecembungan juga tidak awet
terhadap penyusunan tanpa kemonotonan
(\cref{exo:b2:realfun:10}). (ii) Jensen terbalik untuk fungsi
cekung: separuh ketaksamaan klasiknya adalah versi $\ln$ yang
cekung; dan menerapkan bentuk cembungnya pada $\ln$ adalah cara
tercepat membuktikan AM--GM secara \emph{terbalik}. (iii) Kecembungan
titik-tengah semata tidak berarti kecembungan --- diperlukan
kekontinuan (atau sekadar keterbatasan) (\cref{exo:b2:realfun:8});
adapun contoh penyangkal yang patologis hidup di luar aksioma buku ini.
(iv) Fungsi cembung pada interval \emph{\omterm{def:b2:metric:topology}{terbuka}} bersifat \omterm{def:b2:metric:continuity}{kontinu},
bahkan \omterm{def:b2:metric:continuity}{Lipschitz} secara lokal (\cref{exo:b2:realfun:12}); sedangkan di
ujungnya, tak ada yang gratis. (v) Turunan menuruti Darboux tetapi
tidak harus \omterm{def:b2:metric:continuity}{kontinu} (\cref{ex:b2:realfun:oscillation}):
``$f'$ tak punya lompatan'' tak pernah berarti ``$f'$ \omterm{def:b2:metric:continuity}{kontinu}''.
\end{remark}

\section{Sifat Darboux}

\begin{theorem}[Darboux]\label{thm:b2:realfun:darboux}
Misalkan $f$ dapat diturunkan pada sebuah interval $I$. Maka $f'$
mengambil setiap nilai di antara dua nilainya --- walaupun $f'$ tidak
harus \omterm{def:b2:metric:continuity}{kontinu}.\index{teorema Darboux}
\end{theorem}

\begin{proof}
Misalkan $a < b$ di $I$ dan $v$ tegas di antara $f'(a)$ dan $f'(b)$,
katakanlah $f'(a) < v < f'(b)$. Fungsi $g(x) = f(x) - vx$ dapat
diturunkan dengan $g'(a) < 0 < g'(b)$: jadi minimumnya pada
$\intcc{a}{b}$ (yang tercapai berkat kekontinuan pada himpunan \omterm{def:b2:metric:compact}{kompak})
tidak di $a$ (sebab tepat sesudah $a$, $g$ turun di bawah $g(a)$) maupun
di $b$ (sebab tepat sebelum $b$, $g$ berada di bawah $g(b)$): jadi ia di
bagian dalam, dan di sana $g'(c) = 0$, yakni $f'(c) = v$. (Ini dulu
latihan berbintang Tahun ke-1; tempatnya di dalam teori ada di sini.)
\end{proof}

\begin{example}[Fungsi mana yang merupakan turunan?]\label{ex:b2:realfun:whichderivatives}
Teorema Darboux adalah mesin ketakadaan. Fungsi lantai
$\lfloor x\rfloor$ bukan turunan fungsi apa pun pada
$\R$: sebab ia mengambil nilai $0$ dan $1$ tetapi melewatkan $\frac12$
pada $\intcc{0}{1}$, dan \cref{thm:b2:realfun:darboux} melarang hal itu
bagi turunan. Vonis yang sama menimpa setiap fungsi yang berlompatan
--- fungsi tanda, Heaviside, semua fungsi tangga --- betapa pun polos
rupanya; ``antiturunannya'' ($\abs x$ untuk fungsi tanda, dan
seterusnya) hanya ada di luar lompatannya lalu bersimpul di sana dengan
sebuah sudut. Bandingkan: turunan $f'$ yang liar tak \omterm{def:b2:metric:continuity}{kontinu} pada
\cref{ex:b2:realfun:oscillation} \emph{memang} sebuah turunan --- sebab
ketakkontinuannya berupa ayunan, yang dimaklumi Darboux. Batas
antara kedua perilaku itu persis akibat ``tanpa lompatan'' di bawah ini.
\end{example}

\begin{corollary}\label{cor:b2:realfun:nojumps}
Sebuah turunan tak punya ketakkontinuan berupa lompatan: jika $f'(a^-)$
dan $f'(a^+)$ ada, maka keduanya sama dengan $f'(a)$. Ketakkontinuan
sebuah turunan selalu bertipe ayunan (seperti turunan
$x^2\sin\frac1x$ di $0$, jilid Tahun ke-1).
\end{corollary}

\begin{proof}
Jika $f'(a^+) = \lim_{x\to a^+} f'(x)$ ada dan berbeda dari
$f'(a)$, maka nilai yang tegas di antara keduanya akan dilewatkan $f'$
pada persekitaran kanan --- dan itu bertentangan dengan Darboux pada
interval $\intcc{a}{a + h}$. (Cara lain: teorema nilai rata-rata memaksa
$f'(a) = \lim_{h\to0^+} \frac{f(a+h)-f(a)}{h} = f'(a^+)$, sebab hasil
bagi selisihnya adalah nilai $f'$ di suatu titik antara.)
Demikian pula di sebelah kiri.
\end{proof}

\begin{example}[Turunan berayun yang kanonik]\label{ex:b2:realfun:oscillation}
Misalkan $f(x) = x^2\sin\frac1x$ untuk $x \neq 0$ dan $f(0) = 0$. Di
$0$: $\bigl|\frac{f(h) - f(0)}{h}\bigr| = \abs{h\sin\frac1h}
\leq \abs h \to 0$, jadi $f'(0) = 0$ memang ada. Di luar $0$,
\[
f'(x) = 2x\sin\frac1x - \cos\frac1x ,
\]
yang suku pertamanya menuju $0$ sedangkan $\cos\frac1x$ berayun
melewati $\intcc{-1}{1}$ pada setiap interval $\intoo{0}{\delta}$: jadi
limit $f'(0^+)$ tidak ada. Karenanya $f'$ terdefinisi di mana-mana tetapi
tidak \omterm{def:b2:metric:continuity}{kontinu} di $0$ --- dan, persis seperti diramalkan
\cref{cor:b2:realfun:nojumps}, ketakkontinuannya berupa
ayunan, bukan lompatan: sebab pada tiap $\intoo{0}{\delta}$, $f'$ masih
menyapu seinterval penuh di sekitar $0$. Turunan boleh liar, tetapi
hanya dengan cara yang selaras dengan Darboux.
\end{example}

\begin{remark}[Di mana bab ini dipakai]
Kecembungan adalah mesin bagi industri ketaksamaan: soal akhir pekan
bab ini memproduksi Young, Hölder, Minkowski
dan rantai rata-rata pangkat darinya, yang lalu dikonsumsi teori \omterm{def:b2:nvs:norm}{norma}
pada \cref{ch:b2:nvs} dan taksiran integral pada \cref{ch:b2:integration};
sedangkan Jensen muncul kembali di peluang sebagai ketaksamaan momen pada
\cref{ch:b2:randomvar}. Keteraturan monoton kembali pada
\cref{ch:b2:integration} (sebab fungsi monoton terintegralkan) dan,
pada jilid Tahun ke-3, sebagai keterdiferensialan hampir di mana-mana
bagi fungsi monoton --- tempat ``lompatan yang \omterm{def:b2:structures:countable}{terbilang}'' menjadi
langkah pertama teori Lebesgue.
\end{remark}

\section{Latihan}

\begin{exercise}[$\star$]\label{exo:b2:realfun:1}
Tentukan himpunan ketakkontinuan beserta besar lompatannya: $\lfloor x
\rfloor$; $\;x - \lfloor x\rfloor$; $\;\lfloor x \rfloor + \sqrt{x
- \lfloor x\rfloor}$; dan fungsi pada contoh sesudah
\cref{thm:b2:realfun:monotone} yang dibatasi pada bilangan rasional
diadik $r_n$.
\end{exercise}

\begin{exercise}[$\star$]\label{exo:b2:realfun:2}
Buktikan bahwa fungsi naik $f \colon I \to \R$ yang punya sifat nilai
antara (yakni petanya atas sembarang subinterval berupa interval)
pastilah \omterm{def:b2:metric:continuity}{kontinu}.
\end{exercise}

\begin{exercise}[$\star$]\label{exo:b2:realfun:3}
Manakah di antara berikut ini yang cembung pada daerah asalnya? $x
\mapsto x\ln x$ ($x > 0$); $\;x \mapsto \ln(1 + \eu^x)$; $\;x \mapsto
\sqrt{1 + x^2}$; dan $\;x \mapsto x^3$.
\end{exercise}

\begin{exercise}[$\star\star$]\label{exo:b2:realfun:4}
Misalkan $f$ cembung pada $\R$ dan terbatas di atas. Buktikan bahwa $f$
konstan. \emph{(Jika $f(a) \neq f(b)$, ketaksamaan kemiringannya
merambatkan kemiringan tali busur yang tak nol: melampaui titik bernilai
lebih besar, $f$ tumbuh sedikitnya secara linear --- yang bertentangan
dengan keterbatasannya. Tangani kedua tanda kemiringannya.)} Turunkan
bahwa fungsi cembung pada $\R$ yang berasimtot di kedua ujungnya adalah afin.
\end{exercise}

\begin{exercise}[$\star\star$]\label{exo:b2:realfun:5}
Misalkan $f$ dapat diturunkan pada $I$ dengan $f'$ yang monoton.
Buktikan bahwa $f'$ \omterm{def:b2:metric:continuity}{kontinu} \emph{(padukan
\cref{thm:b2:realfun:monotone} dan
\cref{cor:b2:realfun:nojumps})}.
\end{exercise}

\begin{exercise}[$\star\star$]\label{exo:b2:realfun:6}
(Rata-rata geometri sebagai limit) Untuk $x_i$ yang positif dan bobot
$\lambda_i$ berjumlah $1$, buktikan
\[
\lim_{p \to 0^+} \Bigl(\sum_i \lambda_i x_i^p\Bigr)^{1/p}
= \prod_i x_i^{\lambda_i} ,
\]
lewat $x_i^p = \eu^{p\ln x_i} = 1 + p\ln x_i + O(p^2)$, lalu turunkan
ketaksamaan AM--GM berbobot dari
\cref{ex:b2:realfun:powermeans}.
\end{exercise}

\begin{exercise}[$\star\star$]\label{exo:b2:realfun:7}
(Ketaksamaan entropi) Dengan memakai kecembungan tegas $t \mapsto t\ln
t$, buktikan bahwa untuk $p_i, q_i$ positif dengan $\sum p_i = \sum q_i = 1$:
\[
\sum_i p_i \ln\frac{p_i}{q_i} \geq 0 ,
\]
dengan kesamaan bila dan hanya bila $p = q$. \emph{(Tulis ruas kirinya
sebagai $\sum q_i\, \varphi\bigl(\frac{p_i}{q_i}\bigr)$ dengan
$\varphi(t) = t\ln t$ lalu terapkan Jensen dengan bobot $q_i$.)}
\end{exercise}

\begin{exercise}[$\star\star\star$]\label{exo:b2:realfun:8}
(Kecembungan titik-tengah) Fungsi $f \colon I \to \R$ disebut
\emph{cembung titik-tengah} apabila selalu berlaku
$f\bigl(\frac{x+y}{2}\bigr) \leq \frac{f(x) + f(y)}{2}$. Buktikan bahwa
fungsi cembung titik-tengah yang \emph{\omterm{def:b2:metric:continuity}{kontinu}} pastilah cembung.
\emph{(Tegakkan ketaksamaan kecembungannya untuk bobot diadik
$\frac{k}{2^m}$ secara induktif pada $m$, lalu lewatkan ke limitnya
dengan memakai kepadatan dan kekontinuannya.)}
\end{exercise}

\begin{exercise}[$\star\star\star$]\label{exo:b2:realfun:9}
Misalkan $f$ cembung pada $\intco{0}{+\infty}$ dengan $f(0) \leq 0$.
Buktikan bahwa $x \mapsto \frac{f(x)}{x}$ naik pada
$\intoo{0}{+\infty}$, lalu turunkan bahwa untuk $f$ cembung dengan $f(0)
= 0$ berlaku $f(x + y) \geq f(x) + f(y)$ untuk $x, y \geq 0$
(yakni kesuperaditifan).
\end{exercise}

\begin{exercise}[$\star$]\label{exo:b2:realfun:10}
Misalkan $f$ cembung pada $I$ dan $g$ cembung serta \emph{naik} pada
sebuah interval yang memuat $f(I)$. Buktikan bahwa $g \circ f$ cembung,
lalu tunjukkan lewat contoh penyangkal bahwa kemonotonan $g$ tak dapat
dihilangkan.
\end{exercise}

\begin{exercise}[$\star\star$]\label{exo:b2:realfun:11}
(Hermite--Hadamard) Misalkan $f$ cembung dan \omterm{def:b2:metric:continuity}{kontinu} pada
$\intcc{a}{b}$. Buktikan
\[
f\Bigl(\frac{a+b}{2}\Bigr) \;\leq\; \frac{1}{b -
a}\int_a^b f(t)\,\dd t \;\leq\; \frac{f(a) + f(b)}{2} .
\]
\emph{(Kiri: integralkan sebuah garis penyangga di titik tengahnya.
Kanan: batasi $f$ oleh tali busurnya.)}
\end{exercise}

\begin{exercise}[$\star\star\star$]\label{exo:b2:realfun:12}
Buktikan bahwa fungsi cembung pada interval \emph{\omterm{def:b2:metric:topology}{terbuka}} $I$ bersifat
\omterm{def:b2:metric:continuity}{Lipschitz} secara lokal: untuk setiap ruas $\intcc{a}{b} \subseteq I$
dan marjin $\delta > 0$ dengan $\intcc{a - \delta}{b + \delta}
\subseteq I$, pembatasan $f$ pada $\intcc{a}{b}$ bersifat
\omterm{def:b2:metric:continuity}{Lipschitz}, dengan konstanta
$\max\Bigl(\bigl|\frac{f(a) - f(a - \delta)}{\delta}\bigr|,
\bigl|\frac{f(b + \delta) - f(b)}{\delta}\bigr|\Bigr)$
\emph{(jebaklah tiap kemiringan tali busurnya di antara kedua nilai itu
lewat ketaksamaan kemiringan)}.
\end{exercise}

\section{Soal: Kotak Perkakas Kecembungan}

Satu definisi --- tali busur di atas grafiknya --- menghasilkan
seluruh kotak perkakas ketaksamaan klasik. Soal akhir pekan ini
membangunnya menurut urutan logisnya: kriteria kecembungan dan Jensen
yang tegas, lalu Young, Hölder dan Minkowski (yakni akta kelahiran
norma-$p$), lalu rantai rata-rata pangkat yang \omterm{def:b2:metric:complete}{lengkap} dari minimum ke
maksimum, dan dua panen mahkota --- ketaksamaan Carleman, dan
Hölder yang dibaca sebagai dualitas. Semuanya dibuktikan; tak ada yang
diimpor.

\begin{problem}[{Soal akhir pekan --- Young, Hölder, Minkowski,
dan rantai rata-rata pangkat}]
\label{pb:b2:realfun:1}
Di sepanjang soal ini, $p, q > 1$ merupakan \emph{eksponen sekawan}:
$\frac1p + \frac1q = 1$; vektornya $a = (a_1, \dots, a_n) \in \R^n$;
dan bobot $\lambda_i > 0$ memenuhi $\sum_i\lambda_i = 1$.

\medskip
\noindent\textbf{Bagian I --- Kriteria dan Jensen yang tegas.}
\begin{enumerate}
  \item Misalkan $f$ dapat diturunkan pada sebuah interval $I$.
        Buktikan bahwa $f$ cembung bila dan hanya bila $f'$ naik
        \emph{(satu arahnya lewat pelewatan ke limit pada
        ketaksamaan kemiringan \cref{lem:b2:realfun:slopes};
        arah satunya lewat teorema nilai rata-rata)}. Lalu turunkan
        kriteria $C^2$-nya, yakni $f'' \geq 0$.
  \item Andaikan $f''> 0$ pada $I$. Buktikan bahwa $f$ cembung
        \emph{tegas} (yakni ketaksamaan tegas untuk $x \neq y$
        dan $\lambda \in \intoo01$), dan bahwa fungsi cembung tegas
        memenuhi ketaksamaan Jensen
        (\cref{thm:b2:realfun:convexreg}~(3)) dengan kesamaan
        \emph{hanya} ketika semua $x_i$ berimpit.
  \item Sahkan bahan mentah kotak perkakasnya: $-\ln$ cembung
        tegas pada $\intoo{0}{+\infty}$; $t \mapsto t^r$ cembung
        tegas di sana untuk $r > 1$ dan cekung tegas
        untuk $0 < r < 1$; dan $\exp$ cembung tegas pada $\R$.
  \item (Ketaksamaan Young) Untuk $a, b \geq 0$, buktikan
        \[
        ab \;\leq\; \frac{a^p}{p} + \frac{b^q}{q},
        \]
        dengan kesamaan bila dan hanya bila $a^p = b^q$ \emph{(terapkan
        kecekungan $\ln$ pada kedua titik $a^p, b^q$
        dengan bobot $\frac1p, \frac1q$)}.
  \item Turunkan kembali AM--GM berbobot dalam satu baris dari
        kecekungan $\ln$:
        \[
        \prod_i x_i^{\lambda_i} \leq \sum_i\lambda_ix_i \qquad
        (x_i > 0),
        \]
        beserta kasus kesamaannya; lalu bandingkan dengan jalur limit
        pada \cref{exo:b2:realfun:6}.
\end{enumerate}

\medskip
\noindent\textbf{Bagian II --- Hölder dan Minkowski.} Tulis
$\norm{a}_p = \bigl(\sum_i \abs{a_i}^p\bigr)^{1/p}$ dan
$\norm{a}_\infty = \max_i\abs{a_i}$.
\begin{enumerate}[resume]
  \item (Hölder) Buktikan
        \[
        \sum_{i=1}^{n}\abs{a_ib_i} \;\leq\;
        \norm a_p\,\norm b_q ,
        \]
        dengan kesamaan bila dan hanya bila vektor $(\abs{a_i}^p)$ dan
        $(\abs{b_i}^q)$ sebanding \emph{(normalkan
        $\norm a_p = \norm b_q = 1$ lalu terapkan Young suku demi suku)}.
  \item Kenali kasus khususnya: $p = q = 2$
        (yakni Cauchy--Schwarz), dan pasangan ujungnya $(p, q) = (1,
        \infty)$: nyatakan dan buktikan $\sum\abs{a_ib_i} \leq
        \norm a_1\norm b_\infty$.
  \item (Minkowski) Untuk $p \geq 1$, buktikan
        \[
        \norm{a + b}_p \leq \norm a_p + \norm b_p
        \]
        \emph{(tulis $\abs{a_i + b_i}^p \leq \abs{a_i +
        b_i}^{p-1}(\abs{a_i} + \abs{b_i})$ lalu terapkan Hölder
        pada tiap hasil kalinya)}. Simpulkan: $\norm\cdot_p$ merupakan
        \omterm{def:b2:nvs:norm}{norma} pada $\R^n$ untuk setiap $p \in \intco{1}{+\infty}$,
        sehingga melengkapi gambaran \cref{ch:b2:nvs}.
  \item Versi integralnya: untuk $f, g$ yang \omterm{def:b2:metric:continuity}{kontinu} pada
        $\intcc{a}{b}$, nyatakan dan buktikan Hölder dan Minkowski
        bagi $\norm f_p = \bigl(\int_a^b\abs
        f^p\bigr)^{1/p}$ \emph{(dengan bukti yang sama, memakai
        kepositifan tegas integralnya untuk membahas
        kesamaannya)}.
  \item Buktikan kemonotonan $\norm a_q \leq \norm a_p$ untuk $1
        \leq p \leq q$, limit $\norm a_p \to
        \norm a_\infty$ ketika $p \to \infty$, dan perbandingan
        sebaliknya dengan konstanta yang tajam:
        \[
        \norm a_p \leq n^{\frac1p - \frac1q}\,\norm a_q
        \]
        \emph{(lewat Hölder terhadap vektor konstan)}. Kenali
        vektor yang mencapai tiap kesamaannya.
  \item (Interpolasi) Untuk $1 \leq p < r < q$ dan $\theta \in
        \intoo01$ dengan $\frac1r = \frac\theta p +
        \frac{1-\theta}q$, buktikan
        \[
        \norm a_r \leq \norm a_p^{\theta}\,
        \norm a_q^{1-\theta}
        \]
        \emph{(terapkan Hölder dengan eksponen $\frac{p}{\theta
        r}$ dan $\frac{q}{(1-\theta)r}$ pada $\abs{a_i}^{\theta
        r}\abs{a_i}^{(1-\theta)r}$)}.
\end{enumerate}

\medskip
\noindent\textbf{Bagian III --- Rantai rata-rata pangkat, selengkapnya.}
Untuk $p \neq 0$ tetapkan $M_p = \bigl(\sum_i\lambda_i
x_i^p\bigr)^{1/p}$ (dengan $x_i > 0$), dan $M_0 = \prod_i
x_i^{\lambda_i}$.
\begin{enumerate}[resume]
  \item Buktikan bahwa $p \mapsto M_p$ naik pada seluruh
        $\R^*$: tangani $p < q < 0$ lewat kesamaan kebalikan
        $M_{-p}(x) = M_p(1/x)^{-1}$, lalu jembatani lewat $0$ dengan
        menunjukkan $M_p \leq M_0 \leq M_q$ untuk $p < 0 < q$
        \emph{(terapkan kecekungan $\ln$ pada $x_i^q$, dan
        ketaksamaan terbaliknya untuk eksponen negatif)}.
  \item Buktikan limitnya, yakni $M_p \to \max_i x_i$ ketika $p \to
        +\infty$ dan $M_p \to \min_i x_i$ ketika $p \to -\infty$.
  \item Tuliskan rantai $\min \leq \mathrm{HM} \leq
        \mathrm{GM} \leq \mathrm{AM} \leq \mathrm{QM} \leq
        \max$ untuk bobot yang sama, lalu buktikan akibat
        klasiknya: untuk $a_1, \dots, a_n$ yang positif,
        \[
        \Bigl(\sum_i a_i\Bigr)\Bigl(\sum_i\frac1{a_i}\Bigr)
        \geq n^2 .
        \]
  \item Kaitkan rata-rata dengan \omterm{def:b2:nvs:norm}{norma}: untuk bobot sama $\lambda_i =
        \frac1n$ berlaku $M_p(x) = n^{-1/p}\norm x_p$. Damaikanlah
        kedua kemonotonannya --- rata-rata \emph{naik} terhadap $p$
        sedangkan \omterm{def:b2:nvs:norm}{norma} \emph{turun} (pertanyaan 10) --- dalam satu
        kalimat tentang faktor $n^{-1/p}$.
  \item Tentukan kasus kesamaan di sepanjang seluruh rantai pada
        pertanyaan 14 (dengan bobot positif): kesamaan di mana pun
        memaksa semua $x_i$ sama --- jadi kecembungan tegas berbuah.
\end{enumerate}

\medskip
\noindent\textbf{Bagian IV --- Panennya.}
\begin{enumerate}[resume]
  \item (Young dengan kenop) Untuk $a, b \geq 0$ dan $\varepsilon >
        0$, buktikan
        \[
        ab \leq \varepsilon\,\frac{a^p}{p} +
        \varepsilon^{-q/p}\,\frac{b^q}{q},
        \]
        beserta kasus kuda bebannya $ab \leq \varepsilon a^2 +
        \frac{b^2}{4\varepsilon}$: yakni siasat penyerapan yang dipakai
        di sepanjang analisis.
  \item (Menuju Carleman) Misalkan $c_k = \frac{(k+1)^k}{k^{k-1}}$.
        Buktikan kesamaan teleskopnya $\prod_{k=1}^{n}c_k =
        (n+1)^n$, lalu turunkan, lewat AM--GM yang diterapkan pada
        bilangan $c_ka_k$,
        \[
        (a_1a_2\cdots a_n)^{1/n} \leq
        \frac{1}{n(n+1)}\sum_{k=1}^{n} c_k a_k
        \qquad (a_k > 0).
        \]
  \item (Ketaksamaan Carleman) Jumlahkan atas $n$, tukarkan
        urutan penjumlahannya (lewat \omterm{def:b2:series:summable}{keluarga terjumlahkan} yang positif,
        \cref{thm:b2:series:fubini}), lalu pakai $\sum_{n \geq
        k}\frac{1}{n(n+1)} = \frac1k$ dan $c_k/k = \bigl(1 +
        \frac1k\bigr)^k < \eu$ untuk menyimpulkan: untuk setiap
        $\sum a_k$ yang konvergen dan bersuku positif,
        \[
        \sum_{n=1}^{\infty}(a_1a_2\cdots a_n)^{1/n}
        \;\leq\; \eu\sum_{k=1}^{\infty}a_k .
        \]
  \item Untuk $f$ yang \omterm{def:b2:metric:continuity}{kontinu} dan positif pada $\intcc{0}{1}$, buktikan
        \[
        \Bigl(\int_0^1 f\Bigr)\Bigl(\int_0^1\frac1f\Bigr)
        \geq 1,
        \]
        dengan kesamaan bila dan hanya bila $f$ konstan
        \emph{(lewat Cauchy--Schwarz pada $\sqrt f\cdot\frac1{\sqrt
        f}$)}.
  \item (Geometri bolanya) Dengan memakai kasus kesamaan
        Minkowski, tunjukkan bahwa untuk $1 < p < \infty$ permukaan
        bola satuan $\norm\cdot_p$ tak memuat satu pun ruas (yakni
        \omterm{def:b2:nvs:norm}{normanya} cembung tegas dalam arti
        \cref{pb:b2:nvs:1}), sedangkan untuk $p = 1$ dan $p =
        \infty$ ia memuatnya: tunjukkanlah potongan datarnya.
\end{enumerate}

\medskip
\noindent\textbf{Bagian V --- Dualitas dan rangkuman.}
\begin{enumerate}[resume]
  \item (Hölder sebagai dualitas) Buktikan bahwa untuk setiap $a \in
        \R^n$,
        \[
        \norm a_p = \max_{\norm b_q \leq 1}\ \sum_i a_ib_i ,
        \]
        beserta $b$ yang memaksimumkannya secara gamblang. (Jadi
        norma-$p$ adalah dual norma-$q$ --- yakni benih berdimensi
        hingga bagi dualitas $L^p$.)
  \item (Momen) Misalkan $X$ peubah acak yang mengambil berhingga
        banyak nilai positif $x_i$ dengan peluang
        $\lambda_i$. Nyatakan ulang pertanyaan 12 sebagai: $r \mapsto
        \E[X^r]^{1/r}$ naik --- yakni ketaksamaan momen
        (Lyapunov), yang akan dipakai ulang pada
        \cref{ch:b2:randomvar}.
  \item Selesaikan dengan perkakas yang bernama, masing-masing dua
        baris: (i) untuk $a, b, c$ yang positif: $a^3 + b^3 + c^3 \geq
        \frac{(a+b+c)^3}{9}$; (ii) untuk $x_1, \dots, x_n$ yang positif:
        $\bigl(\sum_i\sqrt{x_i}\bigr)^2 \leq n\sum_i x_i$.
  \item (Rangkuman) Gambarkan silsilahnya dalam lima kalimat: dari
        definisi tali busur ke lema kemiringan; dari kemiringan ke
        garis penyangga ke Jensen; dari kecekungan $\ln$ ke Young ke
        Hölder ke Minkowski ke norma-$p$; dari Jensen ke rantai
        rata-rata pangkat ke momen; dan dari AM--GM ke Carleman.
        Sebutkan puncaknya (Hölder--Minkowski; Carleman), lalu
        nyatakan ke mana kotak perkakas ini menuju: yakni ruang $L^p$
        pada jilid Tahun ke-3, yang aksiomanya persis pertanyaan 6 dan 8.

\end{enumerate}
\end{problem}
